\documentclass{article}
\usepackage[final]{neurips_2024}
\usepackage[utf8]{inputenc}
\usepackage[T1]{fontenc}
\usepackage{hyperref}
\usepackage{url}
\usepackage{booktabs}
\usepackage{amsfonts}
\usepackage{amsmath}
\usepackage{nicefrac}
\usepackage{microtype}
\usepackage{xcolor}
\usepackage{graphicx}
\usepackage{float}
\usepackage{tcolorbox}
\tcbset{
  finding/.style={
    colback=gray!8,
    colframe=black,
    fonttitle=\bfseries,
    title=Takeaway~#1,
    sharp corners,
    boxrule=0.4pt,
    left=6pt, right=6pt, top=4pt, bottom=4pt
  }
}

\title{Adaptive Interaction Graphs for Particle Simulation}

\author{%
  Aiden Zhou \\
  Yale University \\
  \texttt{aiden.zhou@yale.edu}
}

\begin{document}

\maketitle

\begin{abstract}
Learned particle simulators based on graph neural networks achieve strong one-step accuracy, but errors compound over long horizons. An underexplored variable is the interaction graph: existing methods fix its topology via $k$-nearest neighbors or a static radius rule, regardless of local model confidence. We propose making this graph adaptive: a per-particle variance head, trained jointly with the acceleration head under a
heteroscedastic Gaussian NLL loss, drives a trajectory in which high-uncertainty particles receive an expanded neighborhood. This is done at no extra inference cost by using the previous step's uncertainty estimate. A key discovery is that the variance head learns a meaningful notion of uncertainty. Regions of high variance correspond to splash zones, collisions, and free surfaces. When this signal drives graph topology, the resulting \textsc{AdaptGNS} simulator achieves a strict Pareto improvement on \textsc{WaterDrop} and a modest gain on \textsc{Sand}. Given the model's stronger performance on \textsc{WaterDrop} — and perhaps other fluid dynamics tasks — we hypothesize that adaptive graphs are most useful when physical complexity is concentrated in space, a property that \textsc{Sand} does not share.

\end{abstract}



% -----------------------------------------------------------------------
\section{Background} 
% -----------------------------------------------------------------------

\paragraph{Problem statement.}
Learned simulators based on graph neural networks (GNNs) have shown strong one-step accuracy on particle dynamics tasks, but are hurt by compounding error over many steps \cite{gns2020}. A key architectural choice in such systems is the \textbf{interaction graph}, which decides what particles exchange messages at each time step. Existing methods fix this topology using static geometric heuristics, such as a radius rule \cite{gns2020} or a learned edge mask \cite{kipf2018}, regardless of the model's local confidence.

We propose \textsc{AdaptGNS}, an \textbf{uncertainty-adaptive interaction graph} that selectively expands edges per timestep based on the simulator's predicted uncertainty. By providing greater relational capacity where the model is least confident — and leaving the default sparse connectivity in place everywhere else — \textsc{AdaptGNS} aims to make graph network simulators (GNSs) more stable, but also to preserve the speed and efficiency that motivate using a learned simulator in the first place.

\paragraph{Motivation and significance.}
Long-horizon physical simulation is key to a wide range of subjects, such as robotics, molecular dynamics, and world models. Learned simulators can provide much faster inference than classical solvers, 
with reported speedups of up to 5{,}000$\times$ over the Material Point Method (MPM) on granular flow  simulations \cite{gns_joss}. However, this benefit is lost if errors accumulate over many steps. Enhancing consistency without increasing average compute per step would make learned simulators more viable to replace classical solvers.

\paragraph{Learning objective.}
Our task is per-particle acceleration prediction at each timestep, with the predicted values integrated forward to produce a complete simulation trajectory. Errors at each step compound over many steps, and our goal is to reduce this error growth relative to baseline models. A secondary objective trains a \textbf{per-particle uncertainty head} — a scalar variance estimate $\sigma_i$ — that drives the adaptive graph policy at each step. We report two metrics throughout: \textbf{MSE} for standardized accuracy between datasets and models and \textbf{average edges per step} for efficiency. Note that each processor layer of a GNS does $O(|E|)$ work, so compute scales with edge count. 

\paragraph{Summary of results.} \textsc{AdaptGNS} achieves a strict Pareto improvement on \textsc{WaterDrop} (20\% fewer edges and lower MSE at step 200 (MSE@200) than fixed $k=5$), but only a modest gain on \textsc{Sand} (12\% lower MSE@200 than fixed $k=5$, using 10\% fewer edges than $k=10$).\footnote{$k$ is roughly the average number of neighbors per particle at the given connectivity radius $r$. We use this easy notation throughout, but be aware that $r$ is fixed; neighbor count varies with local particle density.} The story becomes clearer once we ablate for the choice of loss function: compared to a fixed GNS \textbf{trained on NLL loss}, the adaptive policy once more achieves a strong gain (21\% fewer edges and 8\% lower MSE@200). We also report the optimized hyperparameters found by grid sweep on \textsc{Sand} and \textsc{WaterDrop} ($\sigma$-percentile $= 70$, radius factor $= 1.267$) and several negative results.

\section{Dataset}

\paragraph{Description.}
We use two distinct particle simulation datasets. \textsc{Sand} is taken from \cite{gns_joss}, hosted on DesignSafe-CI (DOI: \texttt{10.17603/ds2-0phb-dg64}), and generated with the Taichi-MPM solver. \textsc{WaterDrop} is taken from \cite{gns2020} and generated with a Smoothed Particle Hydrodynamics (SPH) solver. Note that \textsc{WaterDrop} trajectories are given as TFRecords, so we provide a script to convert the dataset to a standard \texttt{.npz} format. These two materials were chosen to evaluate adaptive graphs across both fluid and granular flow regimes.

\begin{figure}[H]
    \centering
    \includegraphics[width=0.6\linewidth]{figures/fig_graph_snapshot.pdf}
    \caption{Interaction graph snapshots at Step 100 for \textsc{WaterDrop} (left) and \textsc{Sand} (right).}
    \label{fig:graph_snapshot}
\end{figure}

\paragraph{Graph statistics.} The edges of a standard interaction graph are constructed at each step between particles separated by at most $r = 0.015$. Table~\ref{tab:graph_stats} reports graph statistics taken at Step 100, averaged across trajectories. Both datasets produce sparse graphs with moderate average degree. Notably, both 
datasets have an average clustering coefficient of roughly $0.50$, indicating that 
roughly half of each particle's neighbor-pairs are also mutual neighbors.

\begin{table}[h]
  \caption{Graph statistics at Step 100}
  \label{tab:graph_stats}
  \centering
  \begin{tabular}{lcccccc}
    \toprule
    Dataset & Trajs & Nodes & Edges & Density & Avg Degree & Avg Clustering \\
    \midrule
    \textsc{WaterDrop} & 30  & $567 \pm 200$  & $1550 \pm 583$  
              & $1.08 \times 10^{-2}$ & $5.42 \pm 0.22$ & $0.494 \pm 0.036$ \\
    \textsc{Sand}      & 30 & $1270 \pm 503$ & $4324 \pm 1864$ 
              & $7.35 \times 10^{-3}$ & $6.76 \pm 1.01$ & $0.500 \pm 0.034$ \\
    \bottomrule
  \end{tabular}
\end{table}

The temporal evolution of graph 
statistics reveals varying dynamics: in \textsc{WaterDrop}, 
edge count and average degree rise as the droplet spreads and then taper as particles disperse, whereas in \textsc{Sand} the graph undergoes a sharp slump around Step 50 before recovering to a stable plateau.

\paragraph{Feature analysis.}
\begin{itemize}
    \item \textbf{Node features:} the last 5 velocity vectors, concatenated with a particle type indicator.
    \item \textbf{Edge features:} relative displacement vector and distance between particles.
    \item \textbf{Labels:} per-particle acceleration at the next step.
\end{itemize}

\textbf{Data splits.} We use those provided alongside each dataset. In particular, note that trajectories are independent and data is split at the trajectory level. For both \textsc{Sand} and \textsc{WaterDrop}, we leverage the complete training set and evaluate performance on a fixed subset of 30 test trajectories. Hyperparameter tuning was conducted on the validation set only. 

\begin{table}[h]
  \caption{Data splits per dataset}
  \label{tab:data_splits}
  \centering
  \begin{tabular}{lcccc}
    \toprule
    Dataset & Train trajectories & Validation trajectories & Test trajectories & Steps / trajectory \\
    \midrule
    \textsc{Sand}      & 1{,}000 & 100 & 30 & 320   \\
    \textsc{WaterDrop} & 1{,}000 & 30  & 30 & 1{,}000 \\
    \bottomrule
  \end{tabular}
\end{table}



\section{Related Work}
The idea of representing physical systems as graphs was introduced by \cite{battaglia2016}, who developed the relational inductive bias that subsequent work extends upon. \cite{gns2020} scaled this paradigm to large particle 
systems via radius-based interaction graphs and noise-augmented training. This work forms the 
primary baseline and codebase for our project. A central restriction shared by both 
works is that the interaction graph is fixed by geometric rules, which do not regard local dynamics. \cite{topping2022} provide the theoretical 
grounding for why we care: data from distant nodes is sharply compressed
through bottleneck edges (over-squashing), and this worsens over long periods to cause compounding error.

The most straightforward try to learn the graph structure itself came from \cite{kipf2018}, whose Neural Relational Inference (NRI) derives a discrete interaction graph via a VAE encoder over particle histories. However, NRI produces a graph that is fixed at inference, and incurs significant overhead. More recently, AdaMeshNet \cite{adameshnet2025} proposes 
per-step adaptation for mesh-based GNNs using geometric heuristics, outperforming 
static baselines. Our work diverges from both: 
we target point-cloud particle systems rather than meshes, use no separate encoder, and use the model's own predicted uncertainty $\sigma_i$ as 
the adaptation signal. This leads to a lightweight, general architecture that adapts at every step as a byproduct of the dynamics head itself.

% -----------------------------------------------------------------------
\section{Methodology}
% -----------------------------------------------------------------------

\paragraph{Variance head.}
We extend a standard GNS with a \textbf{variance prediction head}: alongside the acceleration output, the network predicts a scalar uncertainty $\sigma_i$ per particle. This is done via a two-layer MLP on the final node embeddings, using a softplus activation to enforce positivity: \[
\sigma_i^2 = \operatorname{softplus}(f_{\theta}(h_i)).
\]
This change adds fewer than 20,000 parameters relative to the standard GNS, so the architectural overhead is inconsequential. See \texttt{VarianceHead} in
\url{adaptive-gns/gns/graph_network.py}. 

\paragraph{Heteroscedastic NLL loss.} Both acceleration and variance heads are trained jointly under a Gaussian NLL loss:
\begin{equation}
  \mathcal{L} = \frac{1}{N}\sum_{i=1}^{N} \frac{(\hat{a}_i - a_i)^2}{2\sigma_i^2} + \frac{1}{2}\log \sigma_i^2.
\end{equation}

In contrast to MSE, which provides no gradient to $\sigma_i$, NLL ensures that the model is rewarded for concentrating uncertainty where errors are large and punished for inflating it elsewhere. 

\paragraph{Adaptive graph construction loop.} Our core contribution is an inference-time loop that uses predicted uncertainty to modify the graph topology. A naive implementation would run a GNN forward pass, extract
$\sigma_i$, rebuild the graph, and then run a second forward pass to compute
the final acceleration. However, this doubles inference cost. Instead, we use a \textbf{single-pass lagged-$\sigma$} design. At step
$t$, the graph is constructed using the previous step's uncertainty estimates
$\sigma_i^{(t-1)}$. The forward pass at step $t$ then produces both the next
position and a fresh uncertainty estimate $\sigma_i^{(t)}$ for
step $t + 1$. 

Concretely, our adaptive mechanism consists of three operations:
\begin{enumerate}
    \item Construct the standard interaction graph (with radius $r$).
    \item Identify particles whose $\sigma_i$ is above a
    percentile threshold $\tau$.
    \item Add extra edges from a larger radius
    $r_{\text{large}} = \texttt{radius\_factor} \cdot r$ for edges incident to one or more high-uncertainty particles.
\end{enumerate}
See \texttt{predict\_positions\_adaptive} and its helper functions in \url{adaptive-gns/gns/learned_simulator.py} for more information and implementation details.
% -----------------------------------------------------------------------
\section{Experiments and Results}
% -----------------------------------------------------------------------

\subsection{Experimental setup.}

\paragraph{Hyperparameters.}
The adaptive policy has two primary hyperparameters: the uncertainty
threshold $\tau$ and the radius expansion factor \texttt{radius\_factor}. Particles
above the $\tau$-th percentile in predicted variance receive expanded
neighborhoods, whereas other particles keep the default connectivity.

We choose \texttt{radius\_factor} so that the enlarged adaptive radius matches
the radius used by the fixed $k=10$ baseline. As the base $k=5$ graph uses
$r = 0.015$ and the $k=10$ case corresponds to roughly
$r = 0.019$, we set $
  \texttt{radius\_factor}
  =
  \frac{0.019}{0.015}
  \approx 1.267.$
This gives high-$\sigma$ particles in the adaptive model exactly the same connectivity as the strongest fixed-graph baseline.
We verify this choice by tuning $\tau$ and \texttt{radius\_factor} jointly via a grid search over $\tau \in \{65, 70, 75, 80, 90\}$ and \texttt{radius\_factor} $\in \{1.2, 1.267, 1.5, 2\}$. Indeed, we found that $\tau = 70$, \texttt{radius\_factor} $= 1.267$ is the optimal configuration on the validation set.
\paragraph{Baselines and Ablations.} We compare \textsc{AdaptGNS} to five models on both \textsc{WaterDrop} and \textsc{Sand}. The naive comparison is to a \textbf{per-particle MLP}, which measures the value of relational inductive bias. Far more reasonable models are a \textbf{fixed-graph GNS with $k = 5$} (or $r = 0.015$), which serves as the primary comparison, and a \textbf{fixed-graph GNS with $k = 10$} (or $r = 0.019$), which gives an upper bound on the benefit of denser interaction graphs. 

We also conduct two ablation studies. \textbf{GNS ($k$-mid)} uses an intermediate radius ($r=0.0176$) matched to \textsc{AdaptGNS}'s mean edge count, isolating whether our gains come from adaptive topology or just from having more edges on average.\footnote{Note that this ablation was not conducted for \textsc{WaterDrop}. On \textsc{Sand}, $k_\text{mid}$ was constructed by interpolating in $r^2$ between the $k = 5$ and $k = 10$ edge counts to match \textsc{AdaptGNS}'s budget, resulting in $r \approx 0.0176$. On \textsc{WaterDrop}, \textsc{AdaptGNS} already uses \textbf{fewer} edges than fixed $k = 5$, so no equivalent comparison exists.}
\textbf{\textsc{AdaptGNS} fixed-graph uses NLL} but with fixed $r=0.015$, isolating the contribution of the adaptive loop from the changed loss function.

\paragraph{Computational details.} Everything was run on Yale's Misha cluster, and scheduled via SLURM jobs in \texttt{jobs/}. Each 500,000-step training run takes roughly 6 hours on an Nvidia A100 GPU, and evaluating a model on the test set of \textsc{WaterDrop} or \textsc{Sand} takes at most 30 minutes.

\subsection{Quantitative Results}

Our first question is whether the variance head reveals physical structure. If $\sigma_i$ is just noise, a graph policy that depends on it is doubtful to succeed. Fortunately, we come to a positive conclusion.

\begin{figure}[H]
    \centering
    \includegraphics[width=0.9\linewidth]{sigma_traj0_strip_1_10_50.pdf}
    \caption{Per-particle uncertainty $\sigma_i$ at steps 1, 10, and 50 of a \textsc{Sand} trajectory. High-$\sigma$ particles are concentrated in impact zones, free surfaces, and regions of rapid velocity change.}
    \label{fig:sigma_spatial}
\end{figure}

\begin{tcolorbox}[finding=1]
\textbf{As shown in Figure 2, the network correctly infers what ``uncertainty'' means in the context of physical simulation, purely as a byproduct of the learning objective. }This is supported by our calculations in Table 3. On \textsc{Sand}, Bins 0–8 are tightly calibrated, with $|\Delta| \leq 0.06$. The total ECE is large only because the model severely underestimates its error on the hardest particles (Bin 9). Results are weaker on \textsc{WaterDrop}, but conform to the same trend. As our adaptive methodology thresholds at a percentile rank, this does not pose a huge obstacle.
\end{tcolorbox}

\begin{table}[H]
  \caption{Per-decile calibration on \textsc{Sand} and \textsc{WaterDrop}.}
  \label{tab:calibration_sand_wd}
  \centering
  \footnotesize
  \begin{tabular}{c|ccc|ccc}
    \toprule
    & \multicolumn{3}{c|}{\textsc{Sand}} & \multicolumn{3}{c}{\textsc{WaterDrop}} \\
    Bin & mean SE & $2\hat{\sigma}^2$ & $|\Delta|$ & mean SE & $2\hat{\sigma}^2$ & $|\Delta|$ \\
    \midrule
    0 & 0.023 & 0.009 & 0.014 & 0.543 & 0.003 & 0.540 \\
    1 & 0.025 & 0.014 & 0.011 & 0.540 & 0.004 & 0.537 \\
    2 & 0.024 & 0.018 & 0.006 & 0.131 & 0.009 & 0.121 \\
    3 & 0.026 & 0.023 & 0.003 & 0.081 & 0.015 & 0.066 \\
    4 & 0.033 & 0.029 & 0.004 & 0.208 & 0.022 & 0.186 \\
    5 & 0.046 & 0.038 & 0.008 & 0.495 & 0.032 & 0.463 \\
    6 & 0.067 & 0.051 & 0.016 & 1.386 & 0.048 & 1.339 \\
    7 & 0.141 & 0.084 & 0.057 & 3.167 & 0.065 & 3.101 \\
    8 & 0.560 & 0.195 & 0.365 & 5.739 & 0.094 & 5.645 \\
    9 & 15.482 & 1.882 & 13.599 & 17.878 & 0.288 & 17.590 \\
    \midrule
    \multicolumn{1}{l}{Total ECE} & & & 1.408 & & & 2.959 \\
    \bottomrule
  \end{tabular}
\end{table}

Next, we evaluate \textsc{AdaptGNS} on \textsc{Sand} and \textsc{WaterDrop}. See Table 4 for a summary of results and Figure 3 for more nuanced dynamics.

\paragraph{\textsc{Sand}.} Comparing \textsc{AdaptGNS} to its fixed-graph counterpart (using NLL and $r = 0.015$), the adaptive policy \textbf{reduces edges by 21\% and
lowers MSE@200 by 8\% (0.0417 vs. 0.0382)}.
This is the cleanest comparison: it isolates the value of placing edges
purposely rather than uniformly.
Second, comparing \textsc{AdaptGNS} to $k_\text{mid}$ ($r = 0.0176$) at matched edge counts
(${\sim}23$,000 edges each), \textbf{the model using MSE wins (0.0355 vs.\ 0.0382)}. The benefit of MSE as an easier learning objective is made clear by looking at the NLL model at $r = 0.0176$. Although it also uses a larger fixed radius, its MSE@200 is much worse (0.0422).

\paragraph{\textsc{WaterDrop}.}
The result here is an unequivocal advance: \textsc{AdaptGNS} \textbf{uses 20\%
fewer edges than $k = 5$ (7,821 vs.\ 9,741) and also achieves lower MSE@200
(0.0412 vs.\ 0.0414).}
Even compared to $k = 10$, the tradeoff is 36\% fewer edges for a 6\% relative MSE@200
increase. In contrast to \textsc{Sand}, the decrease in performance caused by using NLL instead of MSE is close to zero.

\begin{table}[H]
  \caption{Results on \textsc{Sand} (314 steps) and \textsc{WaterDrop} (994 steps).}
  \label{tab:rollouts_sand_wd_mse200}
  \centering
  \footnotesize
  \setlength{\tabcolsep}{5pt}
  \begin{minipage}[t]{0.48\linewidth}
    \centering
    \textbf{\textsc{Sand}}\\[4pt]
    \begin{tabular}{@{}lrr@{}}
      \toprule
      Model & Edges/step & MSE@200 \\
      \midrule
      MLP & N/A & 0.121 \\
      $k =5$ ($r = 0.015$) & 18{,}682 & 0.0434 \\
      $k_{\mathrm{mid}}$ ($r = 0.0176$) & 23{,}498 & 0.0355 \\
      $k=10$ ($r = 0.019$) & 25{,}516 & \textbf{0.0285} \\
      NLL ($r = 0.015$) & 28{,}968 & 0.0417 \\
      NLL ($r = 0.0176$) & 28{,}768 & 0.0422 \\
      \textbf{AdaptGNS} & \textbf{22{,}958} & 0.0382 \\
      \bottomrule
    \end{tabular}
  \end{minipage}\hfill
  \begin{minipage}[t]{0.48\linewidth}
    \centering
    \textbf{\textsc{WaterDrop}}\\[4pt]
    \begin{tabular}{@{}lrr@{}}
      \toprule
      Model & Edges/step & MSE@200 \\
      \midrule
      MLP & N/A & 0.207 \\
      $k=5$ ($r = 0.015$) & 9{,}741 & 0.0414 \\
      $k=10$ ($r = 0.019$) & 12{,}212 & \textbf{0.0388} \\
      NLL ($r = 0.015$) & 7{,}875 & 0.0417 \\
      \textbf{AdaptGNS} & \textbf{7{,}821} & 0.0412 \\
      \bottomrule
    \end{tabular}
  \end{minipage}
\end{table}

\begin{figure}[H]
    \centering
    \includegraphics[width=\linewidth]{edges_and_mse_sand_waterdrop1.png}
    \caption{Per-step dynamics (edge count and MSE) on \textsc{Sand} and \textsc{WaterDrop}.}
    \label{fig:sigma_spatial}
\end{figure}

\begin{tcolorbox}[finding=2]
\textbf{\textsc{AdaptGNS} strictly beats the same NLL model without the adaptive graph mechanism on both datasets.} On \textsc{WaterDrop}, it also achieves a strong advance over every fixed-graph comparison. Due to its harder learning objective, \textsc{AdaptGNS} does not yet defeat the best MSE-trained models on \textsc{Sand}.
\end{tcolorbox}

\subsection{Analysis and Discussion}

\paragraph{Hypotheses revisited.} To \textbf{RQ1} (does adaptive expansion reduce MSE at increasingly long horizons?) and \textbf{RQ2} (does the variance head identify physical structure?), \textbf{we answer yes!} On both \textsc{WaterDrop} and \textsc{Sand}, we observe a significant gain over the base NLL model when the adaptive mechanism is ablated. And as shown in Figure 2 and Table 3, the variance head learns a relevant concept of physical uncertainty. \textbf{RQ3} (match $k = 10$ accuracy with fewer edges?) is true on \textsc{WaterDrop} but only partly on \textsc{Sand}, where the $k_\text{mid}$ ablation shows that a standard GNS (using MSE) at the same edge budget wins on raw accuracy. 

\paragraph{Why \textsc{WaterDrop} and not \textsc{Sand}?}
The performance gap between datasets can be attributed to two physical properties. First, \textbf{graph density evolves differently}. Most granular trajectories end up packing particles into an increasingly dense heap. A droplet impact does the opposite, dispersing particles and keeping the graph sparse for the rest of the trajectory. Second, \textbf{uncertainty is more heterogeneous in \textsc{WaterDrop}}: a splash creates local complexity but is mostly surrounded by quiet, giving the variance head plenty of contrastive signal. In \textsc{Sand}, nearly every particle undergoes the same dynamic, as was previously observed in our exploratory data analysis. 

\section{Conclusion}

In this paper, we explored making the interaction graph of a learned particle simulator \textbf{adaptive} — sparse where the model is confident, denser where it is not — without paying the cost of two GNN computations per step. We contribute a \textbf{single-pass adaptive mechanism} driven by a \textbf{per-particle variance head}, which is trained jointly with the acceleration head under \textbf{NLL loss}. We observe that $\sigma_i$ corresponds to real physical uncertainty, and that \textsc{AdaptGNS}, driven by this signal, achieves notable results on the \textsc{WaterDrop} and \textsc{Sand} datasets. The contrast between datasets suggests that adaptive interaction graphs are most useful for modeling contexts such as fractures, impacts, and turbulent interfaces, where uncertainty is inherently heterogeneous in space and where compounding error is hardest to control by uniform compute scaling.

This conclusion comes with several caveats. For one, switching from MSE to NLL costs around $30\%$ relative MSE@200 on \textsc{Sand}, and the adaptive mechanism reduces but does not erase this gap. Second, the \textbf{single-pass design reacts one step late}. Hence, during the most rapid transitions, $\sigma_i$ may not reflect the current step's uncertainty. Lastly, \textbf{absolute calibration is poor} on both datasets. This is not too serious in our case, but would impact any downstream task that used the actual variance values. Future work should target a learned next-step uncertainty predictor (solving both lag and calibration concerns) and a learnable per-particle expansion radius. We also intend to test \textsc{AdaptGNS} on a broader range of physical regimes (e.g.\ \textsc{Goop}, \textsc{MultiMaterial}) and on longer-horizon tasks. 

\newpage
\bibliographystyle{plainnat}
\bibliography{references}

\end{document}
